\section{Related Work and Positioning}

Units, time and coordinates are represented in several ways across astronomical
software. Astropy integrates quantities, time and coordinate objects in Python,
with attribute-aware transformations between frames \cite{astropy2022}.
SOFA and ERFA provide standard numerical routines for IAU fundamental astronomy,
with conventions expressed through documented arguments and reference
algorithms \cite{sofa,erfa}. SPICE couples geometry to body identifiers,
reference frames, epochs and kernel data \cite{acton1996spice,acton2018spice}.
Orekit organizes flight-dynamics calculations around dates, frames, states and
physical models \cite{orekit2016,orekitSoftware}. These designs demonstrate the value of
scientific representations alongside numerical algorithms.

Rust-native projects bring related design choices to compiled scientific
software. ANISE combines SPICE-inspired geometry, frame metadata and ephemeris
data through Rust and Python interfaces \cite{anise}. Nyx composes spacecraft
states, dynamics, propagation and estimation for mission analysis \cite{nyx}.
Brahe combines Rust and Python satellite-dynamics interfaces, epochs,
coordinate conversions and configurable Earth-orientation data \cite{brahe}.
Hifitime supplies precise epochs, durations and explicit time-scale identifiers
\cite{hifitime}. The \texttt{uom} library represents dimensional quantities
statically and converts their unit representations \cite{uom}.

Astropy carries unit and scale metadata in runtime objects and checks
compatibility during execution; frame attributes include epoch and observer
location. ANISE also enforces scientific context: its frame descriptors carry
origin and orientation identifiers, with fallible kernel-graph traversal
\cite{aniseFrames}. Hifitime stores a time-scale value in its epoch, whereas
\texttt{uom} encodes dimensions and unit systems statically, normalizing
quantities to base units. These are substantive semantic mechanisms, not an
absence of safety. Orekit uses dates, frame trees and model providers; its
standard numerical arguments follow SI conventions. SOFA/ERFA expose convention-
specific scalar and array arguments; SPICE validates names, kernels and coverage
at runtime. Model and data selection remain scientific decisions in each design.

Table~\ref{tab:positioning} compares the architectural boundaries most relevant
to Siderust. Our quantity layer retains unit markers in public types; our time
layer parameterizes epochs by scale and encoding; and our geometry layer
parameterizes positions by frame and center. Siderust connects these generic
representations through an astronomy-specific transformation and model layer.
The same abstractions therefore serve a terrestrial observing workflow, a lunar
surface calculation or a numerical dynamics component. Applications can also
adopt the lower-level libraries independently.

\begin{table*}[!tbp]
\centering
\caption{Architectural approaches to scientific context. Entries describe public
interface choices rather than numerical rankings.}
\label{tab:positioning}
\footnotesize
\begin{tabularx}{\textwidth}{@{}lYYY@{}}
\toprule
System & Scientific representations & Composition mechanism & Relevant design emphasis \\
\midrule
Astropy & Quantities, time and coordinate objects & Dynamic units and frame-attribute checks & Integrated Python astronomy workflows \\
SOFA/ERFA & Convention-specific numerical inputs & Standardized routines and argument conventions & Fundamental-astronomy reference algorithms \\
SPICE & Bodies, frames, epochs and kernels & Metadata-driven kernel geometry & Model-backed planetary and mission geometry \\
Orekit & Dates, frames and spacecraft states & Object and model interfaces & Flight-dynamics and estimation workflows \\
ANISE & Frame metadata, Cartesian states and ephemerides & Rust domain objects and kernel data & Navigation and ephemeris composition \\
Nyx & Trajectory states, Hifitime epochs and dynamical models & Propagation and estimation interfaces & Astrodynamics and mission analysis \\
Brahe & Epochs, numerical states and frame conversions & Rust/Python operations with EOP context & Satellite analysis and propagation \\
Hifitime & Epochs, durations and scale identifiers & Dedicated temporal operations & Reusable precision-time component \\
\texttt{uom} & Dimension-parameterized quantities & Static dimensional arithmetic & Reusable unit-system abstraction \\
Siderust & Unit, scale, frame and center types & Generic scientific layers plus model contexts & Composable astronomical computation \\
\bottomrule
\end{tabularx}
\end{table*}

These projects establish the broader design space in which Siderust operates.
Its distinguishing emphasis is the composition of independently reusable
scientific layers through type-visible astronomical interfaces. We evaluate
that implementation using reference operations with aligned models, frames
and observables, so that numerical comparisons address the intended scientific
computation.

\subsection{A concrete context-preserving workflow}

The Earth-from-Moon calculation in Section~\ref{sec:lunar-example} provides a
common workflow with a finite target, epoch, frame, center and units. Its
independent reference uses the actual SpiceyPy interfaces
\texttt{spkpos("EARTH", et, "J2000", "NONE", "MOON")} and % chktex 36
\texttt{pxform("J2000", "IAU\_MOON", et)}, followed by observer subtraction % chktex 36
and a local projection. CSPICE specifies kilometres and TDB seconds from J2000;
body, frame and correction identifiers are argument values checked against the
loaded kernels \cite{spkposAPI}. These calls make scientific choices explicit
and support a broad runtime frame graph.

In the Siderust listing, \texttt{Position<Selenocentric, EclipticMeanJ2000,
Kilometer>} carries center, frame and unit in its type, and the epoch is
\texttt{JulianDate<TT>}. An explicit sign reversal constructs Earth's
selenocentric position from the Moon's geocentric vector. The surface operation
then selects the frame transformation and preserves the Moon observer identity.
Site longitude, latitude and height, epoch value, loaded BSP and model context
remain runtime values. The compiler rejects incompatible position arithmetic;
fallible calls check site admissibility and kernel access. It cannot establish
that the sign reversal represents the intended target, that two supplied epochs
coincide, or that a selected orientation is adequate for an observation.

Both implementations retain scientific context through different mechanisms.
The distinguishing contribution here is carrying selected contracts across
independently reusable quantity, geometry and time components, then testing their
composition in a domain application. The comparison concerns interface structure;
productivity, error frequency and comparative workflow speed were not measured.
